\section{工业化时代的语言模型}

本讲随后把视角从学术动机转向产业现实：语言模型已经高度工业化了。
课程用 GPT-4、xAI 集群、Stargate 之类的例子来说明，前沿模型的训练成本、参数规模和算力投入已经大到普通研究者无法复现。
而且最关键的是，这些模型的构建细节并不公开。即使是公开的技术报告，很多时候也只会给出很有限的信息。

这意味着两件事：
\begin{itemize}
    \item 前沿模型本身已经超出了课堂的直接可达范围。
    \item 但是这不妨碍我们构建“小模型”来学习机制，只是必须意识到它们未必在规模上完全代表前沿系统。
\end{itemize}

讲师用两个例子说明“规模改变一切”：
\begin{enumerate}
    \item 在小模型里，attention 和 MLP 的 FLOPs 可能差不多；但放大到大模型后，MLP 往往成为主导项。
    \item 某些能力会在规模上“涌现”，在小规模时几乎看不出来。
\end{enumerate}

这两点共同指向一个结论：只在小规模上优化，不一定能优化到真正重要的地方。

\begin{warningbox}{More is different}
规模不是简单的线性放大。
很多结构、瓶颈和能力只会在更大尺度上显现出来。
所以本课程既要教机制，也要教“规模意识”。
\end{warningbox}

